\chapter{Manajemen Risiko Enterprise dan Risiko TI}

\section*{Tujuan Pembelajaran}
\begin{itemize}
  \item Menerapkan proses ISO 31000 untuk mengidentifikasi, menganalisis, dan mengevaluasi risiko terkait TI.
  \item Merumuskan pernyataan \textit{risk appetite} yang dapat dipakai untuk mengambil keputusan.
  \item Membandingkan penilaian risiko kualitatif dan kuantitatif beserta keterbatasan masing-masing.
\end{itemize}

\section{Pendahuluan}

% TODO
% Cakupan: organisasi punya daftar risiko yang panjang, tetapi daftar itu tidak pernah mengubah keputusan apa pun.

\section{Konsep Dasar Risiko}

% TODO
% Cakupan: risiko sebagai efek ketidakpastian terhadap tujuan; ancaman dan peluang; risiko inheren dan risiko residual; peran kontrol.

\section{ISO 31000:2018: Prinsip, Kerangka, dan Proses}

% TODO
% Cakupan: prinsip manajemen risiko; kerangka (kepemimpinan, integrasi, desain, implementasi, evaluasi, perbaikan); proses (konteks, identifikasi, analisis, evaluasi, penanganan, pemantauan, komunikasi); status revisi ISO 31000 yang sedang berjalan di ISO/TC 262 (periksa ulang).

\section{COSO ERM dan Risk Appetite}

% TODO
% Cakupan: lima komponen dan 20 prinsip COSO ERM 2017; risk appetite, risk tolerance, dan kapasitas risiko; ciri pernyataan risk appetite yang dapat dipakai direksi.
% Kasus Medikara: contoh pernyataan risk appetite untuk keselamatan pasien, kerahasiaan data, kelangsungan layanan, dan kepatuhan.

\section{Risiko Keamanan Informasi menurut ISO/IEC 27005:2022}

% TODO
% Cakupan: pendekatan berbasis peristiwa dan berbasis aset; kriteria penerimaan risiko; kaitan dengan klausul penilaian dan penanganan risiko pada ISO/IEC 27001 (dirinci di Bab 6).

\section{Risk Register dan Penanganan Risiko}

% TODO
% Cakupan: struktur risk register; menulis skenario risiko yang baik (sebab, peristiwa, akibat); opsi penanganan (menghindari, mengurangi, membagi, menerima); pemilik risiko dan risiko sisa.
% Kasus Medikara: skenario kebocoran lewat vendor, kontrol akses lemah, gangguan integrasi BPJS Kesehatan, dan laporan manajemen berbasis spreadsheet.

\section{Matriks Risiko dan Pengukuran Kuantitatif}

% TODO
% Cakupan: kelemahan matriks risiko (skor yang menipu, kategori kasar); gambaran pengukuran kuantitatif seperti Open FAIR dan simulasi Monte Carlo; kapan pendekatan kualitatif sudah cukup; pengantar KRI (dirinci di Bab 13).
% Pendalaman teknis: simulasi Monte Carlo di spreadsheet dan loss exceedance curve mengikuti Hubbard dan Seiersen (2023).

\section{Aktivitas: Risk Appetite dan Risk Register Kondisi Saat Ini}

% TODO
% Tujuan aktivitas: menetapkan batas risiko yang dapat diterima dan memotret risiko utama organisasi saat ini.
% Peran dalam tim: analis risiko memimpin; anggota berlatar akuntansi atau keuangan menilai dampak finansial; analis teknologi menilai kemungkinan dan kontrol yang ada; analis regulasi menautkan risiko ke register kewajiban dari Bab 3.
% Langkah: rumuskan risk appetite untuk tiga sampai empat kategori; identifikasi 10 sampai 15 skenario; nilai dampak dan kemungkinan; catat kontrol yang ada; tetapkan pemilik dan opsi penanganan awal.
% Luaran dan tonggak capstone: pernyataan risk appetite dan risk register kondisi saat ini; masuk C2.
% Berkas sources yang direncanakan: sources/chapter05/templat-risk-appetite.md, sources/chapter05/risk-register.xlsx, sources/chapter05/simulasi-monte-carlo.xlsx

\section{Ringkasan}

% TODO

\section{Latihan}

\begin{enumerate}
  \item Ubah pernyataan ``Medikara tidak menoleransi kebocoran data'' menjadi pernyataan \textit{risk appetite} yang realistis dan dapat diukur.
  \item Dua risiko sama-sama mendapat skor 12 pada matriks 5 x 5, tetapi hanya satu yang berdampak pada keselamatan pasien. Jelaskan kelemahan matriks risiko dari contoh ini.
  \item Tentukan opsi penanganan yang paling tepat untuk risiko ketergantungan pada vendor aplikasi Medikara+, beserta alasan dan risiko sisanya.
\end{enumerate}
